% GENERATED FILE. Edit canonical lessons, then run npm run generate:books.
% canonical-source-hash: fnv1a64:ee6148708d1805d9
% canonical-lessons: RU-C46-breakfast, RU-C46-lunch, RU-C46-dinner, RU-C46-porridge, RU-C46-pie

\chapter{Breakfast, Lunch, Dinner}
\label{ch:ru-breakfast-lunch-dinner}
\hlchaptermodality{46}

\begin{chapteropening}
\textbf{By the end of this chapter:} \emph{I can name breakfast, lunch and dinner, porridge and a pie.}

Complete the last lesson of chapter 46: i can name breakfast, lunch and dinner, porridge and a pie.
\end{chapteropening}

\section[závtrak]{\ru{завтрак} (závtrak) — breakfast}

\label{lesson:RU-C46-breakfast}



\begin{quote}
Before the new one: say the Russian for a cucumber, then the Russian for a carrot.
\end{quote}

\subsection*{You'll want to know: \ru{завтрак}}
\textbf{\ru{завтрак}} (\emph{závtrak}) — \textquotedblleft{}breakfast\textquotedblright{}.

\begin{grammarlens}[title={What you've built}]
One more word for this chapter.
\end{grammarlens}

\subsection*{Your turn}
\begin{itemize}
\raggedright
  \item \emph{Say it:} \emph{závtrak}
  \item \emph{Say it:} \emph{závtrak}, once more
  \item \emph{Say it:} say \emph{závtrak}
  \item \emph{Recall:} say the Russian for a cucumber, then the Russian for a carrot, then say \emph{závtrak} again
\end{itemize}

\begin{tcolorbox}[breakable,colback=teal!4,colframe=teal!35!black,title={Before you move on}]
What does \ru{завтрак} mean? (\textquotedblleft{}Breakfast\textquotedblright{}.) Say it once more.
\end{tcolorbox}
\section[obéd]{\ru{обед} (obéd) — lunch, the midday meal}

\label{lesson:RU-C46-lunch}



\begin{quote}
Before the new one: say the Russian for a carrot, then the Russian for breakfast.
\end{quote}

\subsection*{You'll want to know: \ru{обед}}
\textbf{\ru{обед}} (\emph{obéd}) — \textquotedblleft{}lunch, the midday meal\textquotedblright{}.

\begin{grammarlens}[title={What you've built}]
One more word for this chapter.
\end{grammarlens}

\subsection*{Your turn}
\begin{itemize}
\raggedright
  \item \emph{Say it:} \emph{obéd}
  \item \emph{Say it:} \emph{obéd}, once more
  \item \emph{Say it:} say \emph{obéd}
  \item \emph{Recall:} say the Russian for a carrot, then the Russian for breakfast, then say \emph{obéd} again
\end{itemize}

\begin{tcolorbox}[breakable,colback=teal!4,colframe=teal!35!black,title={Before you move on}]
What does \ru{обед} mean? (\textquotedblleft{}Lunch, the midday meal\textquotedblright{}.) Say it once more.
\end{tcolorbox}
\section[úzhin]{\ru{ужин} (úzhin) — dinner, supper}

\label{lesson:RU-C46-dinner}



\begin{quote}
Before the new one: say the Russian for breakfast, then the Russian for lunch.
\end{quote}

\subsection*{You'll want to know: \ru{ужин}}
\textbf{\ru{ужин}} (\emph{úzhin}) — \textquotedblleft{}dinner, supper\textquotedblright{}.

\begin{grammarlens}[title={What you've built}]
One more word for this chapter.
\end{grammarlens}

\subsection*{Your turn}
\begin{itemize}
\raggedright
  \item \emph{Say it:} \emph{úzhin}
  \item \emph{Say it:} \emph{úzhin}, once more
  \item \emph{Say it:} say \emph{úzhin}
  \item \emph{Recall:} say the Russian for breakfast, then the Russian for lunch, then say \emph{úzhin} again
\end{itemize}

\begin{tcolorbox}[breakable,colback=teal!4,colframe=teal!35!black,title={Before you move on}]
What does \ru{ужин} mean? (\textquotedblleft{}Dinner, supper\textquotedblright{}.) Say it once more.
\end{tcolorbox}
\section[kásha]{\ru{каша} (kásha) — porridge}

\label{lesson:RU-C46-porridge}



\begin{quote}
Before the new one: say the Russian for lunch, then the Russian for dinner.
\end{quote}

\subsection*{You'll want to know: \ru{каша}}
\textbf{\ru{каша}} (\emph{kásha}) — \textquotedblleft{}porridge\textquotedblright{}.

\begin{grammarlens}[title={What you've built}]
One more word for this chapter.
\end{grammarlens}

\subsection*{Your turn}
\begin{itemize}
\raggedright
  \item \emph{Say it:} \emph{kásha}
  \item \emph{Say it:} \emph{kásha}, once more
  \item \emph{Say it:} say \emph{kásha}
  \item \emph{Recall:} say the Russian for lunch, then the Russian for dinner, then say \emph{kásha} again
\end{itemize}

\begin{tcolorbox}[breakable,colback=teal!4,colframe=teal!35!black,title={Before you move on}]
What does \ru{каша} mean? (\textquotedblleft{}Porridge\textquotedblright{}.) Say it once more.
\end{tcolorbox}
\section[piróg]{\ru{пирог} (piróg) — a pie}

\label{lesson:RU-C46-pie}



\begin{quote}
Before the new one: say the Russian for dinner, then the Russian for porridge.
\end{quote}

\subsection*{You'll want to know: \ru{пирог}}
\textbf{\ru{пирог}} (\emph{piróg}) — \textquotedblleft{}a pie\textquotedblright{}.

\begin{grammarlens}[title={What you've built}]
One more word for this chapter.
\end{grammarlens}

\subsection*{Your turn}
\begin{itemize}
\raggedright
  \item \emph{Say it:} \emph{piróg}
  \item \emph{Say it:} \emph{piróg}, once more
  \item \emph{Say it:} say \emph{piróg}
  \item \emph{Recall:} say the Russian for dinner, then the Russian for porridge, then say \emph{piróg} again
\end{itemize}

\begin{tcolorbox}[breakable,colback=teal!4,colframe=teal!35!black,title={Before you move on}]
What does \ru{пирог} mean? (\textquotedblleft{}A pie\textquotedblright{}.) Say it once more.
\end{tcolorbox}
